\section{AIME I 1992}
        \begin{enumerate}
        	\item (\textit{AIME I 1992}) Find the sum of all positive rational numbers that are less than 10 and that have denominator 30 when written in lowest terms.


 
 \textbf{Answer: }400



\textbf{Solution 1}

There are 8 \href{https://artofproblemsolving.com/wiki/index.php?title=Fraction}{fractions} which fit the conditions between 0 and 1: $\frac{1}{30},\frac{7}{30},\frac{11}{30},\frac{13}{30},\frac{17}{30},\frac{19}{30},\frac{23}{30},\frac{29}{30}$

 Their sum is 4. Note that there are also 8 terms between 1 and 2 which we can obtain by adding 1 to each of our first 8 terms. For example, $1+\frac{19}{30}=\frac{49}{30}.$ Following this pattern, our answer is $4(10)+8(1+2+3+\cdots+9)=\boxed{400}.$




\textbf{Solution 2}

By \href{https://artofproblemsolving.com/wiki/index.php?title=Euler\%27s\_Totient\_Function}{Euler's Totient Function}, there are $8$ numbers that are relatively prime to $30$, less than $30$. Note that they come in pairs $(m,30-m)$ which result in sums of $1$; thus the sum of the smallest $8$ rational numbers satisfying this is $\frac12\cdot8\cdot1=4$. Now refer to solution 1.

 



\textbf{Solution 3}

Note that if $x$ is a solution, then $(300-x)$ is a solution. We know that $\phi(300) = 80.$ Therefore the answer is $\frac{80}{2} \cdot\frac{300}{30} = \boxed{400}.$



	\item (\textit{AIME I 1992}) A positive integer is called ascending if, in its decimal representation, there are at least two digits and each digit is less than any digit to its right. 
How many ascending positive integers are there?


 
 \textbf{Answer: }502



\textbf{Solution 1}

Note that an ascending number is exactly determined by its \href{https://artofproblemsolving.com/wiki/index.php?title=Digit}{digits}: for any \href{https://artofproblemsolving.com/wiki/index.php?title=Set}{set} of digits (not including 0, since the only position for 0 is at the leftmost end of the number, i.e. a leading 0), there is exactly one ascending number with those digits.

 So, there are nine digits that may be used: $1,2,3,4,5,6,7,8,9.$  Note that each digit may be present or may not be present. Hence, there are $2^9=512$ potential ascending numbers, one for each \href{https://artofproblemsolving.com/wiki/index.php?title=Subset}{subset} of $\{1, 2, 3, 4, 5, 6, 7, 8, 9\}$.

 However, we've counted one-digit numbers and the \href{https://artofproblemsolving.com/wiki/index.php?title=Empty\_set}{empty set}, so we must subtract them off to get our answer, $512-10=\boxed{502}.$



	\item (\textit{AIME I 1992}) A tennis player computes her win ratio by dividing the number of matches she has won by the total number of matches she has played. At the start of a weekend, her win ratio is exactly $0.500$. During the weekend, she plays four matches, winning three and losing one. At the end of the weekend, her win ratio is greater than $.503$. What's the largest number of matches she could've won before the weekend began?


 
 \textbf{Answer: }164



\textbf{Solution 1}

Let $n$ be the number of matches won, so that $\frac{n}{2n}=\frac{1}{2}$, and $\frac{n+3}{2n+4}>\frac{503}{1000}$. 

 Cross \href{https://artofproblemsolving.com/wiki/index.php?title=Multiply}{multiplying}, $1000n+3000>1006n+2012$, so $n<\frac{988}{6}=164 \dfrac {4}{6}=164 \dfrac{2}{3}$. Thus, the answer is $\boxed{164}$.

 



\textbf{Solution 2 (same as solution 1)}

Let $n$ be the number of matches she won before the weekend began. Since her win ratio started at exactly .$500 = \frac{1}{2},$ she must have played exactly $2n$ games total before the weekend began. After the weekend, she would have won $n+3$ games out of $2n+4$ total. Therefore, her win ratio would be $(n+3)/(2n+4).$ This means that
 \[\frac{n+3}{2n+4} > .503 = \frac{503}{1000}.\]Cross-multiplying, we get $1000(n+3) > 503(2n+4),$ which is equivalent to $n < \frac{988}{6} = 164.\overline{6}.$ Since $n$ must be an integer, the largest possible value for $n$ is $\boxed{164}.$



	\item (\textit{AIME I 1992}) In Pascal's Triangle, each entry is the sum of the two entries above it. The first few rows of the triangle are shown below.


 
 \[\begin{array}{c@{\hspace{8em}} c@{\hspace{6pt}}c@{\hspace{6pt}}c@{\hspace{6pt}}c@{\hspace{4pt}}c@{\hspace{2pt}} c@{\hspace{2pt}}c@{\hspace{2pt}}c@{\hspace{2pt}}c@{\hspace{3pt}}c@{\hspace{6pt}} c@{\hspace{6pt}}c@{\hspace{6pt}}c} \vspace{4pt} \text{Row 0: } &    &    &     &     &    &    & 1 &     &     &    &    &    &  \\\vspace{4pt} \text{Row 1: } &    &    &     &     &    & 1 &    & 1  &     &    &    &    &  \\\vspace{4pt} \text{Row 2: } &    &    &     &     & 1 &    & 2 &     & 1  &    &    &    &  \\\vspace{4pt} \text{Row 3: } &    &    &     &  1 &    & 3 &    & 3  &     & 1 &    &    &  \\\vspace{4pt} \text{Row 4: } &    &    & 1  &     & 4 &    & 6 &     & 4  &    & 1 &    &  \\\vspace{4pt} \text{Row 5: } &    & 1 &     & 5  &    &10&    &10 &     & 5 &    & 1 &  \\\vspace{4pt} \text{Row 6: } & 1 &    & 6  &     &15&    &20&     &15 &    & 6 &    & 1 \end{array}\]In which row of Pascal's Triangle do three consecutive entries occur that are in the ratio $3: 4: 5$? 


 
 \textbf{Answer: }062



\textbf{Solution 1}

Consider what the ratio means. Since we know that they are consecutive terms, we can say
 \[\frac{\dbinom{n}{k-1}}{3} = \frac{\dbinom{n}{k}}{4} = \frac{\dbinom{n}{k+1}}{5}.\]

 Taking the first part, and using our expression for $n$ choose $k$,	
 \[\frac{n!}{3(k-1)!(n-k+1)!} = \frac{n!}{4k!(n-k)!}\] 
 \[\frac{1}{3(k-1)!(n-k+1)!} = \frac{1}{4k!(n-k)!}\] 
 \[\frac{1}{3(n-k+1)} = \frac{1}{4k}\] 
 \[n-k+1 = \frac{4k}{3}\] 
 \[n = \frac{7k}{3} - 1\] 
 \[\frac{3(n+1)}{7} = k\]Then, we can use the second part of the equation.	
 \[\frac{n!}{4k!(n-k)!} = \frac{n!}{5(k+1)!(n-k-1)!}\] 
 \[\frac{1}{4k!(n-k)!} = \frac{1}{5(k+1)!(n-k-1)!}\] 
 \[\frac{1}{4(n-k)} = \frac{1}{5(k+1)}\] 
 \[\frac{4(n-k)}{5} = k+1\] 
 \[\frac{4n}{5}-\frac{4k}{5} = k+1\] 
 \[\frac{4n}{5} = \frac{9k}{5} +1.\]Since we know $k = \frac{3(n+1)}{7}$ we can plug this in, giving us	
 \[\frac{4n}{5} = \frac{9\left(\frac{3(n+1)}{7}\right)}{5} +1\] 
 \[4n = 9\left(\frac{3(n+1)}{7}\right)+5\] 
 \[7(4n - 5) = 27n+27\] 
 \[28n - 35 = 27n+27\] 
 \[n = 62\]We can also evaluate for $k$, and find that $k = \frac{3(62+1)}{7} = 27.$ Since we want $n$, however, our final answer is $\boxed{062.}$ ~$\LaTeX$ by ciceronii

 



\textbf{Solution 2}

Call the row $x=t+k$, and the position of the terms $t-1, t, t+1$. Call the middle term in the ratio $N = \dbinom{t+k}{t} = \frac{(t+k)!}{k!t!}$. The first term is $N \frac{t}{k+1}$, and the final term is $N \frac{k}{t+1}$. Because we have the ratio $3:4:5$,

 $\frac{t}{k+1} = \frac{3}{4}$ and $\frac{k}{t+1} = \frac{5}{4}$.

 $4t = 3k+3$ and $4k= 5t+5$
$4t-3k=3$ $5t-4k=-5$ 

 Solve the equations to get $t= 27, k=35$ and $x = t+k = \boxed{062}$.

 -Solution and LaTeX by jackshi2006, variables and algebra simplified by oinava

 


	\item (\textit{AIME I 1992}) Let $S^{}_{}$ be the set of all rational numbers $r^{}_{}$, $0^{}_{}<r<1$, that have a repeating decimal expansion in the form $0.abcabcabc\ldots=0.\overline{abc}$, where the digits $a^{}_{}$, $b^{}_{}$, and $c^{}_{}$ are not necessarily distinct. To write the elements of $S^{}_{}$ as fractions in lowest terms, how many different numerators are required?


 
 \textbf{Answer: }660



\textbf{Solution}

We consider the method in which repeating decimals are normally converted to fractions with an example:

 $x=0.\overline{176}$
$\Rightarrow 1000x=176.\overline{176}$
$\Rightarrow 999x=1000x-x=176$
$\Rightarrow x=\frac{176}{999}$

 Thus, let $x=0.\overline{abc}$
$\Rightarrow 1000x=abc.\overline{abc}$
$\Rightarrow 999x=1000x-x=abc$
$\Rightarrow x=\frac{abc}{999}$

 If $abc$ is not divisible by $3$ or $37$, then this is in lowest terms. Let us consider the other multiples: $333$ multiples of $3$, $27$ of $37$, and $9$ of both $3$ and $37$, so $999-333-27+9 = 648$, which is the amount that are neither. The $12$ numbers that are multiples of $81$ reduce to multiples of $3$. We have to count these since it will reduce to a multiple of $3$ which we have removed from $999$, but, this cannot be removed since the numerator cannot cancel the $3$.There aren't any numbers which are multiples of $37^2$, so we can't get numerators which are multiples of $37$. Therefore $648 + 12 = \boxed{660}$.

 Note: You can use Euler's totient function to find the numbers relatively prime to 999 which gives 648. 

 


	\item (\textit{AIME I 1992}) For how many pairs of consecutive integers in $\{1000,1001,1002^{}_{},\ldots,2000\}$ is no carrying required when the two integers are added?


 
 \textbf{Answer: }156



\textbf{Solution 1}

For one such pair of consecutive integers, let the smaller integer be $\underline{1ABC},$ where $A,B,$ and $C$ are digits from $0$ through $9.$

 We wish to count the ordered triples $(A,B,C).$ By casework, we consider all possible forms of the larger integer, as shown below.
 \[\begin{array}{c|c|c|c|c|c|c} & & & & & & \\ [-2.5ex] \textbf{Case} & \textbf{Thousands} & \textbf{Hundreds} & \textbf{Tens} & \textbf{Ones} & \textbf{Conditions for No Carrying} & \boldsymbol{\#}\textbf{ of Ordered Triples} \\ [0.5ex] \hline & & & & & & \\ [-2ex] 0\leq C\leq 8 & 1 & A & B & C+1 & 0\leq A,B,C\leq 4 & 5^3 \\ 0\leq B\leq 8; \ C=9 & 1 & A & B+1 & 0 & 0\leq A,B\leq 4; \ C=9 & 5^2 \\ 0\leq A\leq 8; \ B=C=9 & 1 & A+1 & 0 & 0 & 0\leq A\leq 4; \ B=C=9 & 5 \\ A=B=C=9 & 2 & 0 & 0 & 0 & A=B=C=9 & 1 \end{array}\]Together, the answer is $5^3+5^2+5+1=\boxed{156}.$




\textbf{Solution 2}

Consider what carrying means: If carrying is needed to add two numbers with digits $abcd$ and $efgh$, then $h+d\ge 10$ or $c+g\ge 10$ or $b+f\ge 10$. 6.  Consider $c \in \{0, 1, 2, 3, 4\}$. $1abc + 1ab(c+1)$ has no carry if $a, b \in \{0, 1, 2, 3, 4\}$. This gives $5^3=125$ possible solutions. 

 With $c \in \{5, 6, 7, 8\}$, there obviously must be a carry. Consider $c = 9$. $a, b \in \{0, 1, 2, 3, 4\}$ have no carry. This gives $5^2=25$ possible solutions. Considering $b = 9$,  $a \in \{0, 1, 2, 3, 4, 9\}$ have no carry. Thus, the solution is $125 + 25 + 6=\boxed{156}$. 

 



\textbf{Solution 3}

Consider the ordered pair $(1abc , 1abc - 1)$ where $a,b$ and $c$ are digits. We are trying to find all ordered pairs where $(1abc) + (1abc - 1)$ does not require carrying. For the addition to require no carrying, $2a,2b < 10$, so $a,b < 5$ unless $1abc$ ends in $00$, which we will address later.  Clearly, if $c \in \{0, 1, 2, 3, 4 ,5\}$, then adding $(1abc) + (1abc - 1)$ will require no carrying. We have $5$ possibilities for the value of $a$, $5$ for $b$, and $6$ for $c$, giving a total of $(5)(5)(6) = 150$, but we are not done yet.

 We now have to consider the cases where $b,c = 0$, specifically when $1abc \in \{1100, 1200, 1300, 1400, 1500, 1600, 1700, 1800, 1900, 2000\}$. We can see that $1100, 1200, 1300, 1400, 1500$, and $2000$ all work, giving a grand total of $150 + 6 = \boxed{156}$ ordered pairs.

 



\textbf{Solution 4 (Complementary Counting)}

Since there are $2000-1000+1 = 1001$ numbers in the set, this means that there are $1000$ consecutive pairs of integers $(a, b)$ in the set. For a pair to carry over, either the hundreds digits will carry over, the tens digits will carry over, or the ones digits will carry over.

 For the hundreds digits, every consecutive pair in $\{1500, 1501, ..., 2000\}$ will have to carry over, except for the pair $\{1999, 2000\}$ when no carrying occurs, removing $500-1 = 499$ pairs. 

 Now, consider the case when the tens digit is more than 5. Every pair of consecutive integers from  $\{1050,1051, ..., 1100\}$ will have to carry over except for the pair $\{1099, 1100\}$, and since there are 5 of these cases less than $1500$ (namely the integers from $1050-1100$, $1150-1200$, $1250-1300$,  $1350-1400$, and $1450-1500$) this removes $49 \cdot 5 = 245$ pairs from the list. 

 Now, consider the case where the ones digits carry over. In this case, every pair of consecutive integers from $\{1005, 1006, 1007, 1008, 1009\}$ carries over, removing 4 pairs. Since there are 5 possible cases for the tens digit ($0-4$) and 5 possible cases for the hundreds digit, there are $4 \cdot 5 \cdot 5 = 100$ more cases to be excluded.   

 In total, there are $1000 - 499 - 245 - 100 = \boxed{156}$ consecutive pairs.

 


	\item (\textit{AIME I 1992}) Faces $ABC^{}_{}$ and $BCD^{}_{}$ of tetrahedron $ABCD^{}_{}$ meet at an angle of $30^\circ$. The area of face $ABC^{}_{}$ is $120^{}_{}$, the area of face $BCD^{}_{}$ is $80^{}_{}$, and $BC=10^{}_{}$. Find the volume of the tetrahedron.


 
 \textbf{Answer: }320



\textbf{Solution}

Since the area $BCD=80=\frac{1}{2}\cdot10\cdot16$, the perpendicular from $D$ to $BC$ has length $16$.

 The perpendicular from $D$ to $ABC$ is $16 \cdot \sin 30^\circ=8$. Therefore, the volume is $\frac{8\cdot120}{3}=\boxed{320}$.

 



\textbf{Solution 2}

The area of $ABC$ is 120 and $BC$=10, the slant height is 24. Height from $A$ to $BCD$ is $24 \cdot \sin 30^\circ=12$. Since area of $BCD$ is 80, the volume of tetrahedron $ABCD$= $\frac{80\cdot12}{3}=\boxed{320}$.

 


 


	\item (\textit{AIME I 1992}) For any sequence of real numbers $A=(a_1,a_2,a_3,\ldots)$, define $\Delta A^{}_{}$ to be the sequence $(a_2-a_1,a_3-a_2,a_4-a_3,\ldots)$, whose $n^{th}$ term is $a_{n+1}-a_n^{}$. Suppose that all of the terms of the sequence $\Delta(\Delta A^{}_{})$ are $1^{}_{}$, and that $a_{19}=a_{92}^{}=0$. Find $a_1^{}$.


 
 \textbf{Answer: }819



\textbf{Solution 1 (uses calculus)}

Note that the $\Delta$s are reminiscent of differentiation; from the condition $\Delta(\Delta{A}) = 1$, we are led to consider the differential equation
 \[\frac{d^2 A}{dn^2} = 1\]This inspires us to guess a quadratic with leading coefficient 1/2 as the solution; 
 \[a_{n} = \frac{1}{2}(n-19)(n-92)\]as we must have roots at $n = 19$ and $n = 92$.

 Thus, $a_1=\frac{1}{2}(1-19)(1-92)=\boxed{819}$.

 



\textbf{Solution 2}

Let $\Delta^1 A=\Delta A$, and $\Delta^n A=\Delta(\Delta^{(n-1)}A)$.

 Note that in every sequence of $a_i$, $a_n=\binom{n-1}{1}\Delta a_n + \binom{n-1}{2}\Delta^2 a_n +\binom{n-1}{3}\Delta^3 a_n + ...$

 Then $a_n=a_1 +\binom{n-1}{1}\Delta a_1 +\binom{n-1}{2}\Delta^2 a_1 +\binom{n-1}{3}\Delta^3 a_1 + ...$

 Since $\Delta a_1 =a_2 -a_1$, $a_n=a_1 +\binom{n-1}{1}(a_2-a_1) +\binom{n-1}{2}\cdot 1=a_1 +n(a_2-a_1) +\binom{n-1}{2}$
$a_{19}=0=a_1+18(a_2-a_1)+\binom{18}{2}=18a_2-17a_1+153$
$a_{92}=0=a_1+91(a_2-a_1)+\binom{91}{2}=91a_2-90a_1+4095$

 Solving, $a_1=\boxed{819}$.

 



\textbf{Solution 3}

The sequence $\Delta(\Delta A)$ is the second finite difference sequence, and the first $k-1$ terms of this sequence can be computed in terms of the original sequence as shown below.

 $\begin{array}{rcl} a_3+a_1-2a_2&=&1\\ a_4+a_2-2a_3&=&1\\ &\vdots\\ a_k + a_{k-2} - 2a_{k-1} &= &1\\ a_{k+1} + a_{k-1} - 2a_k &=& 1.\\ \end{array}$

 Adding the above $k-1$ equations we find that

 
 \[(a_{k+1} - a_k) = k-1 + (a_2-a_1).\tag{1}\]

 We can sum equation $(1)$ from $k=1$ to $18$, finding
 \[18(a_1-a_2) - a_1 = 153.\tag{2}\]

 We can also sum equation $(1)$ from $k=1$ to $91$, finding
 \[91(a_1-a_2) - a_1 = 4095.\tag{3}\]Finally, $18\cdot (3) - 91\cdot(2)$ gives $a_1=\boxed{819}$.

 Kris17

 



\textbf{Solution 3a (Same Solution but Explained More)}

Start with this (from Solution 3):

 $\begin{array}{rcl} a_3-2a_2+a_1&=&1\\ a_4-2a_3+a_2&=&1\\ a_5-2a_4+a_3&=&1\\ &\vdots\\ a_k-2a_{k-1}+a_{k-2} &= &1\\ \end{array}$

 When every equation is added up, most terms cancel to get $a_k-a_{k-1}-a_2+a_1=k-2$.

 Now add up this series of equations:

 $\begin{array}{rcl} a_{92}-a_{91}-a_2+a_1&=&90\\ a_{91}-a_{90}-a_2+a_1&=&89\\ a_{90}-a_{89}-a_2+a_1&=&88\\ &\vdots\\ a_{20}-a_{19}-a_2+a_1&=&18\\ \end{array}$

 Again, most terms cancel or can be added up to get $a_{92}-a_{19}+73(-a_2+a_1)=\frac{(90+18)(90-18+1)}{2}$ $\therefore -a_2+a_1 = 54$.

 Using the original equation $a_k-a_{k-1}-a_{k-1}+a_{k-2}=1$, for $k=3$ substitute to get $-a_3+a_2=53$.

 Now substitute this equation for $k=4$ and continue until $k=19$:

 $\begin{array}{rcl} -a_3+a_2&=&53\\ -a_4+a_3&=&52\\ -a_5+a_4&=&51\\ &\vdots\\ -a_{19}+a_{18}&=&37\\ \end{array}$

 Terms cancel until $-a_{19}+a_1 = \frac{(54+37)(54-37+1)}{2} = 819$.

 $a_{19} = 0$, so $a_1=\boxed{819}$.

 



\textbf{Solution 4}

Since all terms of $\Delta(\Delta A)$ are 1, we know that $\Delta A$ looks like $(k,k+1,k+2,...)$ for some $k$. This means $A$ looks like $(a_1,a_1+k,a_1+2k+1,a_1+3k+3,a_1+4k+6,...)$. More specifically, $A_n=a_1+k(n-1)+\frac{(n-1)(n-2)}{2}$. Plugging in $a_{19}=a_{92}=0$, we have the following linear system: 
 \[a_1+91k=-4095\] 
 \[a_1+18k=-153\] From this, we can easily find that $k=-54$ and $a_1=\boxed{819}$.Solution by Zeroman

 


 



\textbf{Solution 5}

Since the result of two finite differences of some sequence is a constant sequence, we know that sequence is a quadratic. Furthermore, we know that $f(19) = f(92) = 0$ so the quadratic is $f(x) = a(x-19)(x-92)$ for some constant $a.$ Now we use the conditions that the finite difference is $1$ to find $a.$ We know $f(19) = 0$ and $f(20) = -72a$ and $f(18) = 74a.$ Therefore applying finite differences once yields the sequence $-74a,-72a$ and then applying finite differences one more time yields $2a$ so $a =\frac{1}{2}.$ Therefore $f(1) = 9 \cdot 91 = \boxed{819}.$




\textbf{Solution 6}

Let $a_1=a,a_2=b.$ From the conditions, we have 
 \[a_{n-1}+a_{n+1}=2a_n+1,\] for all $n>1.$ From this, we find that 
 \begin{align*} a_3&=2b+1-a \\ a_4&=3b+3-2a\\ a_5&=4b+6-3a, \end{align*}
or, in general, 
 \[a_n=(n-1)b+\frac{(n-2)(n-1)}{2}-(n-2)a.\] Note: we can easily prove this by induction. Now, substituting $n=19,92,$ we find that 
 \begin{align*} 0=&18b+\frac{17\cdot18}{2}-17a\\ 0=&91b+\frac{90\cdot91}{2}-90a\\ b=&\frac{17a-\frac{17\cdot18}{2}}{18}=\frac{90a-\frac{90\cdot91}{2}}{91}. \end{align*}
Now, cross multiplying, we find that
 \begin{align*} 91\left(17a-\frac{17\cdot18}{2}\right)&=18\left(90a-\frac{90\cdot91}{2}\right)\\ 1547a-\frac{(18\cdot91)\cdot17}{2}&=1620-\frac{(18\cdot91)\cdot90}{2}\\ 73a&=\frac{18\cdot91}{2}\cdot(90-17=73)\\ a=\frac{18\cdot91}{2}=9\cdot91=\boxed{819}. \end{align*}




	\item (\textit{AIME I 1992}) Trapezoid $ABCD^{}_{}$ has sides $AB=92^{}_{}$, $BC=50^{}_{}$, $CD=19^{}_{}$, and $AD=70^{}_{}$, with $AB^{}_{}$ parallel to $CD^{}_{}$. A circle with center $P^{}_{}$ on $AB^{}_{}$ is drawn tangent to $BC^{}_{}$ and $AD^{}_{}$. Given that $AP^{}_{}=\frac mn$, where $m^{}_{}$ and $n^{}_{}$ are relatively prime positive integers, find $m+n^{}_{}$.


 
 \textbf{Answer: }164



\textbf{Solution 1}

Let $AP=x$ so that $PB=92-x.$ Extend $AD, BC$ to meet at $X,$ and note that $XP$ bisects $\angle AXB;$ let it meet $CD$ at $E.$ Using the angle bisector theorem, we let $XB=y(92-x), XA=xy$ for some $y.$ 

 Then $XD=xy-70, XC=y(92-x)-50,$ thus 
 \[\frac{xy-70}{y(92-x)-50} = \frac{XD}{XC} = \frac{ED}{EC}=\frac{AP}{PB} = \frac{x}{92-x},\] which we can rearrange, expand and cancel to get $120x=70\cdot 92,$ hence $AP=x=\frac{161}{3}$. This gives us a final answer of $161+3=\boxed{164}$



	\item (\textit{AIME I 1992}) Consider the region $A^{}_{}$ in the complex plane that consists of all points $z^{}_{}$ such that both $\frac{z^{}_{}}{40}$ and $\frac{40^{}_{}}{\overline{z}}$ have real and imaginary parts between $0^{}_{}$ and $1^{}_{}$, inclusive. What is the integer that is nearest the area of $A^{}_{}$?


 
 \textbf{Answer: }572



\textbf{Solution}

Let $z=a+bi \implies \frac{z}{40}=\frac{a}{40}+\frac{b}{40}i$. Since $0\leq \frac{a}{40},\frac{b}{40}\leq 1$ we have the inequality 
 \[0\leq a,b \leq 40\]which is a square of side length $40$.

 Also, $\frac{40}{\overline{z}}=\frac{40}{a-bi}=\frac{40a}{a^2+b^2}+\frac{40b}{a^2+b^2}i$ so we have $0\leq a,b \leq \frac{a^2+b^2}{40}$, which leads to:
 \[(a-20)^2+b^2\geq 20^2\]
 \[a^2+(b-20)^2\geq 20^2\]

 We graph them:

 
                    \begin{figure}[H]
                        \centering
                        \includegraphics[width=0.5\textwidth]{Solution/AIME_I_1992_10_Solution_Figure_1.png}
                    \end{figure}
                    

 To find the area outside the two circles but inside the square, we want to find the unique area of the two circles. We can do this by adding the area of the two circles and then subtracting out their overlap. There are two methods of finding the area of overlap:

 1. Consider that the area is just the quarter-circle with radius $20$ minus an isosceles right triangle with base length $20$, and then doubled (to consider the entire overlapped area)

 2. Consider that the circles can be converted into polar coordinates, and their equations are $r = 40sin\theta$ and $r = 40cos\theta$. Using calculus with the appropriate bounds, we can compute the overlapped area.

 Using either method, we compute the overlapped area to be $200\pi + 400$, and so the area of the intersection of those three graphs is $40^2-(200\pi + 400) \Rightarrow 1200 - 200\pi \approx 571.68$
$\boxed{572}$



	\item (\textit{AIME I 1992}) Lines $l_1^{}$ and $l_2^{}$ both pass through the origin and make first-quadrant angles of $\frac{\pi}{70}$ and $\frac{\pi}{54}$ radians, respectively, with the positive x-axis. For any line $l^{}_{}$, the transformation $R(l)^{}_{}$ produces another line as follows: $l^{}_{}$ is reflected in $l_1^{}$, and the resulting line is reflected in $l_2^{}$. Let $R^{(1)}(l)=R(l)^{}_{}$ and $R^{(n)}(l)^{}_{}=R\left(R^{(n-1)}(l)\right)$. Given that $l^{}_{}$ is the line $y=\frac{19}{92}x^{}_{}$, find the smallest positive integer $m^{}_{}$ for which $R^{(m)}(l)=l^{}_{}$.


 
 \textbf{Answer: }945



\textbf{Solution}

Let $l$ be a line that makes an angle of $\theta$ with the positive $x$-axis. Let $l'$ be the reflection of $l$ in $l_1$, and let $l''$ be the reflection of $l'$ in $l_2$. 

 The angle between $l$ and $l_1$ is $\theta - \frac{\pi}{70}$, so the angle between $l_1$ and $l'$ must also be $\theta - \frac{\pi}{70}$. Thus,  $l'$ makes an angle of $\frac{\pi}{70}-\left(\theta-\frac{\pi}{70}\right) = \frac{\pi}{35}-\theta$ with the positive $x$-axis.  

 Similarly, since the angle between $l'$ and $l_2$ is $\left(\frac{\pi}{35}-\theta\right)-\frac{\pi}{54}$, the angle between $l''$ and the positive $x$-axis is $\frac{\pi}{54}-\left(\left(\frac{\pi}{35}-\theta\right)-\frac{\pi}{54}\right) = \frac{\pi}{27}-\frac{\pi}{35}+\theta = \frac{8\pi}{945} + \theta$. 

 Thus, $R(l)$ makes an $\frac{8\pi}{945} + \theta$ angle with the positive $x$-axis. So $R^{(n)}(l)$ makes an $\frac{8n\pi}{945} + \theta$ angle with the positive $x$-axis. 

 Therefore, $R^{(m)}(l)=l$ iff $\frac{8m\pi}{945}$ is an integral multiple of $\pi$. Thus, $8m \equiv 0\pmod{945}$. Since $\gcd(8,945)=1$, $m \equiv 0 \pmod{945}$, so the smallest positive integer $m$ is $\boxed{945}$. 

 


	\item (\textit{AIME I 1992}) In a game of \textit{Chomp}, two players alternately take bites from a 5-by-7 grid of unit squares. To take a bite, a player chooses one of the remaining squares, then removes ("eats") all squares in the quadrant defined by the left edge (extended upward) and the lower edge (extended rightward) of the chosen square. For example, the bite determined by the shaded square in the diagram would remove the shaded square and the four squares marked by $\times.$ (The squares with two or more dotted edges have been removed from the original board in previous moves.) 


 
                    \begin{figure}[H]
                        \centering
                        \includegraphics[width=0.5\textwidth]{Problem/AIME_I_1992_12_Problem_Figure_1.png}
                    \end{figure}
                    


 The object of the game is to make one's opponent take the last bite. The diagram shows one of the many subsets of the set of 35 unit squares that can occur during the game of Chomp. How many different subsets are there in all? Include the full board and empty board in your count.


 
 \textbf{Answer: }792



\textbf{Solution 1}

By drawing possible examples of the subset, one can easily see that making one subset is the same as dividing the game board into two parts.

 One can also see that it is the same as finding the shortest route from the upper left hand corner to the lower right hand corner; Such a route would require 5 lengths that go down, and 7 that go across, with the shape on the right "carved" out by the path a possible subset.

 Therefore, the total number of such paths is $\binom{12}{5}=\boxed{792}$




\textbf{Solution 2}

If any square is eaten, the squares to the right of it must also be eaten. Thus, if $a_i$ is the number of squares remaining on row $i$, there is exactly one way the row can be configured (the leftmost $a_i$ squares are uneaten and the ones to the right are eaten.) Additionally, the squares above an eaten square must be eaten, so $\{a_i\}$ is nondecreasing. We can thus write $0 \leq a_1 \leq a_2 \leq \dots \leq a_5 \leq 7$; the number of sequences $\{a_i\}$ satisfying this inequality may be found by Stars and Bars to be $\binom{7+6-1}{6-1} = \boxed{792}$.

 \textbf{Note} This game is similar to an AoPS book.

 \textbf{Links to Branches of Mathematics} This question is one case of the problem of counting \href{https://artofproblemsolving.com/wiki/index.php?title=Order\_ideal&action=edit&redlink=1}{order ideals} or \href{https://artofproblemsolving.com/wiki/index.php?title=Antichain&action=edit&redlink=1}{antichains} in \href{https://artofproblemsolving.com/wiki/index.php?title=Poset}{posets}.  Specifically, it ask for the number of order ideals of the \href{https://artofproblemsolving.com/wiki/index.php?title=Product\_poset&action=edit&redlink=1}{product poset} of the \href{https://artofproblemsolving.com/wiki/index.php?title=Chain}{chain} of length $4$ and the chain of length $6$.

 


	\item (\textit{AIME I 1992}) Triangle $ABC^{}_{}$ has $AB=9^{}_{}$ and $BC: AC=40: 41^{}_{}$. What's the largest area that this triangle can have?


 
 \textbf{Answer: }820



\textbf{Solution 1}

First, consider the \href{https://artofproblemsolving.com/wiki/index.php?title=Triangle}{triangle} in a \href{https://artofproblemsolving.com/wiki/index.php?title=Coordinate\_system}{coordinate system} with \href{https://artofproblemsolving.com/wiki/index.php?title=Vertex}{vertices} at $(0,0)$, $(9,0)$, and $(a,b)$. Applying the \href{https://artofproblemsolving.com/wiki/index.php?title=Distance\_formula}{distance formula}, we see that $\frac{ \sqrt{a^2 + b^2} }{ \sqrt{ (a-9)^2 + b^2 } } = \frac{40}{41}$.

 We want to maximize $b$, the height, with $9$ being the base.

 Simplifying gives $-a^2 -\frac{3200}{9}a +1600 = b^2$.

 To maximize $b$, we want to maximize $b^2$. So if we can write: $b^2=-(a+n)^2+m$, then $m$ is the maximum value of $b^2$ (this follows directly from the \href{https://artofproblemsolving.com/wiki/index.php?title=Trivial\_inequality}{trivial inequality}, because if ${x^2 \ge 0}$ then plugging in $a+n$ for $x$ gives us ${(a+n)^2 \ge 0}$).

 $b^2=-a^2 -\frac{3200}{9}a +1600=-\left(a +\frac{1600}{9}\right)^2 +1600+\left(\frac{1600}{9}\right)^2$.

 $\Rightarrow b\le\sqrt{1600+\left(\frac{1600}{9}\right)^2}=40\sqrt{1+\frac{1600}{81}}=\frac{40}{9}\sqrt{1681}=\frac{40\cdot 41}{9}$.

 Then the area is $9\cdot\frac{1}{2} \cdot \frac{40\cdot 41}{9} = \boxed{820}$.

 



\textbf{Solution 2}

Let the three sides be $9,40x,41x$, so the area is $\frac14\sqrt {(81^2 - 81x^2)(81x^2 - 1)}$ by Heron's formula. By AM-GM, $\sqrt {(81^2 - 81x^2)(81x^2 - 1)}\le\frac {81^2 - 1}2$, and the maximum possible area is $\frac14\cdot\frac {81^2 - 1}2 = \frac18(81 - 1)(81 + 1) = 10\cdot82 = \boxed{820}$. This occurs when $81^2 - 81x^2 = 81x^2 - 1\implies x = \frac {\sqrt {3281}}9$.

 \textbf{Comment} Rigorously, we need to make sure that the equality of the AM-GM inequality is possible to be obtained (in other words, $(81^2 - 81x^2)$ and $(81x^2 - 1)$ can be equal with some value of $x$). MAA is pretty good at generating smooth combinations, so in this case, the AM-GM works; however, always try to double check in math competitions -- the writer of Solution 2 gave us a pretty good example of checking if the AM-GM equality can be obtained.   ~Will\_Dai

 



\textbf{Solution 3}

Let $A, B$ be the endpoints of the side with length $9$. Let $\Gamma$ be the Apollonian Circle of $AB$ with ratio $40:41$; let this intersect $AB$ at $P$ and $Q$, where $P$ is inside $AB$ and $Q$ is outside. Then because $(A, B; P, Q)$ describes a harmonic set, $AP/AQ=BP/BQ\implies \dfrac{\frac{41}{9}}{BQ+9}=\dfrac{\frac{40}{9}}{BQ}\implies BQ=360$. Finally, this means that the radius of $\Gamma$ is $\dfrac{360+\frac{40}{9}}{2}=180+\dfrac{20}{9}$.

 Since the area is maximized when the altitude to $AB$ is maximized, clearly we want the last vertex to be the highest point of $\Gamma$, which just makes the altitude have length $180+\dfrac{20}{9}$. Thus, the area of the triangle is $\dfrac{9\cdot \left(180+\frac{20}{9}\right)}{2}=\boxed{820}$




\textbf{Solution 4 (Involves Basic Calculus)}

We can apply Heron's on this triangle after letting the two sides equal $40x$ and $41x$. Heron's gives

 $\sqrt{\left(\frac{81x+9}{2} \right) \left(\frac{81x-9}{2} \right) \left(\frac{x+9}{2} \right) \left(\frac{-x+9}{2} \right)}$.

 This can be simplified to

 $\frac{9}{4} \cdot \sqrt{(81x^2-1)(81-x^2)}$.

 We can optimize the area of the triangle by finding when the derivative of the expression inside the square root equals 0.

 We have that $-324x^3+13124x=0$, so $x=\frac{\sqrt{3281}}{9}$.

 Plugging this into the expression, we have that the area is $\boxed{820}$.

 $\textbf{-RootThreeOverTwo}$




\textbf{Solution 5}

We can start how we did above in solution 4 to get$\frac{9}{4} * \sqrt{(81x^2-1)(81-x^2)}$.Then, we can notice the inside is a quadratic in terms of $x^2$, which is$-81(x^2)^2+6562x^2-81$. This is maximized when $x^2 = \frac{3281}{81}$.If we plug it into the equation, we get $\frac{9}{4} \cdot \frac{3280}{9} = \boxed{820}$



	\item (\textit{AIME I 1992}) In triangle $ABC^{}_{}$, $A'$, $B'$, and $C'$ are on the sides $BC$, $AC^{}_{}$, and $AB^{}_{}$, respectively. Given that $AA'$, $BB'$, and $CC'$ are concurrent at the point $O^{}_{}$, and that $\frac{AO^{}_{}}{OA'}+\frac{BO}{OB'}+\frac{CO}{OC'}=92$, find $\frac{AO}{OA'}\cdot \frac{BO}{OB'}\cdot \frac{CO}{OC'}$.


 
 \textbf{Answer: }094



\textbf{Solution 1}

Let $K_A=[BOC], K_B=[COA],$ and $K_C=[AOB].$  Due to triangles $BOC$ and $ABC$ having the same base,  
 \[\frac{AO}{OA'}+1=\frac{AA'}{OA'}=\frac{[ABC]}{[BOC]}=\frac{K_A+K_B+K_C}{K_A}.\]  Therefore, we have
 \[\frac{AO}{OA'}=\frac{K_B+K_C}{K_A}\] 
 \[\frac{BO}{OB'}=\frac{K_A+K_C}{K_B}\] 
 \[\frac{CO}{OC'}=\frac{K_A+K_B}{K_C}.\]Thus, we are given 
 \[\frac{K_B+K_C}{K_A}+\frac{K_A+K_C}{K_B}+\frac{K_A+K_B}{K_C}=92.\] Combining and expanding gives 
 \[\frac{K_A^2K_B+K_AK_B^2+K_A^2K_C+K_AK_C^2+K_B^2K_C+K_BK_C^2}{K_AK_BK_C}=92.\] We desire $\frac{(K_B+K_C)(K_C+K_A)(K_A+K_B)}{K_AK_BK_C}.$ Expanding this gives 
 \[\frac{K_A^2K_B+K_AK_B^2+K_A^2K_C+K_AK_C^2+K_B^2K_C+K_BK_C^2}{K_AK_BK_C}+2=\boxed{094}.\]




\textbf{Solution 2 (Mass Points)}

Using \href{https://artofproblemsolving.com/wiki/index.php?title=Mass\_points}{mass points}, let the weights of $A$, $B$, and $C$ be $a$, $b$, and $c$ respectively. 

 Then, the weights of $A'$, $B'$, and $C'$ are $b+c$, $c+a$, and $a+b$ respectively. 

 Thus, $\frac{AO^{}_{}}{OA'} = \frac{b+c}{a}$, $\frac{BO^{}_{}}{OB'} = \frac{c+a}{b}$, and $\frac{CO^{}_{}}{OC'} = \frac{a+b}{c}$.

 Therefore:$\frac{AO}{OA'}\cdot \frac{BO}{OB'}\cdot \frac{CO}{OC'} = \frac{b+c}{a} \cdot \frac{c+a}{b} \cdot \frac{a+b}{c}$ $= \frac{2abc+b^2c+bc^2+c^2a+ca^2+a^2b+ab^2}{abc} =$
$2+\frac{bc(b+c)}{abc}+\frac{ca(c+a)}{abc}+\frac{ab(a+b)}{abc} = 2 + \frac{b+c}{a} + \frac{c+a}{b} + \frac{a+b}{c}$ $= 2 + \frac{AO^{}_{}}{OA'}+\frac{BO}{OB'}+\frac{CO}{OC'} = 2+92 = \boxed{094}$.

 



\textbf{Solution 3}

As in above solutions, find $\sum_{cyc} \frac{y+z}{x}=92$ (where $O=(x:y:z)$ in barycentric coordinates). Now letting $y=z=1$ we get $\frac{2}{x}+2(x+1)=92 \implies x+\frac{1}{x}=45$, and so $\frac{2}{x}(x+1)^2=2(x+\frac{1}{x}+2)=2 \cdot 47 = 94$.

 



\textbf{Solution 4 (Ceva's Theorem)}

A consequence of Ceva's theorem sometimes attributed to Gergonne is that $\frac{AO}{OA'}=\frac{AC'}{C'B}+\frac{AB'}{B'C}$, and similarly for cevians $BB'$ and $CC'$.  Now we apply Gergonne several times and do algebra:

 
 \begin{align*} \frac{AO}{OA'}\frac{BO}{OB'}\frac{CO}{OC'} &= \left(\frac{AB'}{B'C}+\frac{AC'}{C'B}\right) \left(\frac{BC'}{C'A}+\frac{BA'}{A'C}\right)  \left(\frac{CB'}{B'A}+\frac{CA'}{A'B}\right)\\ &=\underbrace{\frac{AB'\cdot CA'\cdot BC'}{B'C\cdot A'B\cdot C'A}}_{\text{Ceva}} + \underbrace{\frac{AC'\cdot BA'\cdot CB'}{C'B\cdot A'C\cdot B'A}}_{\text{Ceva}} + \underbrace{\frac{AB'}{B'C} + \frac{AC'}{C'B}}_{\text{Gergonne}} + \underbrace{\frac{BA'}{A'C} + \frac{BC'}{C'A}}_{\text{Gergonne}} + \underbrace{\frac{CA'}{A'B} + \frac{CB'}{B'A}}_{\text{Gergonne}}\\ &= 1 + 1 + \underbrace{\frac{AO}{OA'} + \frac{BO}{OB'} + \frac{CO}{OC'}}_{92} = \boxed{94} \end{align*}




	\item (\textit{AIME I 1992}) Define a positive integer $n^{}_{}$ to be a factorial tail if there is some positive integer $m^{}_{}$ such that the decimal representation of $m!$ ends with exactly $n$ zeroes. How many positive integers less than $1992$ are not factorial tails?


 
 \textbf{Answer: }396



\textbf{Solution}

Let the number of zeros at the end of $m!$ be $f(m)$. We have $f(m) = \left\lfloor \frac{m}{5} \right\rfloor + \left\lfloor \frac{m}{25} \right\rfloor + \left\lfloor \frac{m}{125} \right\rfloor + \left\lfloor \frac{m}{625} \right\rfloor + \left\lfloor \frac{m}{3125} \right\rfloor + \cdots$. 

 Note that if $m$ is a multiple of $5$, $f(m) = f(m+1) = f(m+2) = f(m+3) = f(m+4)$. 

 Since $f(m) \le \frac{m}{5} + \frac{m}{25} + \frac{m}{125} + \cdots  = \frac{m}{4}$, a value of $m$ such that $f(m) = 1991$ is greater than $7964$. Testing values greater than this yields $f(7975)=1991$.

 There are $\frac{7975}{5} = 1595$ distinct positive integers, $f(m)$, less than $1992$. Thus, there are $1991-1595 = \boxed{396}$ positive integers less than $1992$ that are not factorial tails.

 



\textbf{Solution 2}

After testing various values of $m$ in $f(m)$ of solution 1 to determine $m$ for which $f(m) = 1992$, we find that $m \in \{7980, 7981, 7982, 7983, 7984\}$. WLOG, we select $7980$. Furthermore, note that every time $k$ reaches a multiple of $25$, $k!$ will gain two or more additional factors of $5$ and will thus skip one or more numbers.

 With this logic, we realize that the desired quantity is simply $\left \lfloor \frac{7980}{25} \right \rfloor + \left \lfloor \frac{7980}{125} \right \rfloor \cdots$, where the first term accounts for every time $1$ number is skipped, the second term accounts for each time $2$ numbers are skipped, and so on. Evaluating this gives us $319 + 63 + 12 + 2 = \boxed{396}$. - Spacesam(edited by srisainandan6)

 


\end{enumerate}